% Explicit computation of the K3-fibre specialisation of the stage-one
% holomorphic factorisation algebra on K3 x E.
% Author: Raeez Lorgat.

\providecommand{\kBKM}{\kappa_{\mathrm{BKM}}}
\providecommand{\kch}{\kappa_{\mathrm{ch}}}
\providecommand{\kcat}{\kappa_{\mathrm{cat}}}
\providecommand{\kfib}{\kappa_{\mathrm{fiber}}}
\providecommand{\kHeis}{\kappa_{\mathrm{ch}}^{\mathrm{Heis}}}
\providecommand{\Eone}{E_1}
\providecommand{\Etwo}{E_2}
\providecommand{\Ethree}{E_3}
\providecommand{\PhiFA}{\Phi^{\mathrm{FA}}}
\providecommand{\SpCh}{\mathrm{Sp}^{\mathrm{ch}}}
\providecommand{\ZZ}{\mathbb{Z}}
\providecommand{\QQ}{\mathbb{Q}}
\providecommand{\CC}{\mathbb{C}}
\providecommand{\cC}{\mathcal{C}}
\providecommand{\cF}{\mathcal{F}}
\providecommand{\cO}{\mathcal{O}}
\providecommand{\cM}{\mathcal{M}}
\providecommand{\cH}{\mathcal{H}}
\providecommand{\fg}{\mathfrak{g}}
\providecommand{\fh}{\mathfrak{h}}
\providecommand{\fn}{\mathfrak{n}}
\providecommand{\Muk}{\mathrm{Muk}}
\providecommand{\Obs}{\mathrm{Obs}}
\providecommand{\BPS}{\mathrm{BPS}}
\providecommand{\Uchir}{U^{\mathrm{ch}}}
\providecommand{\CoHA}{\mathrm{CoHA}}
\providecommand{\BM}{\mathrm{BM}}
\providecommand{\BorLift}{\mathrm{BorLift}}
\providecommand{\dbar}{\bar\partial}

\section{The specialisation \texorpdfstring{$\SpCh_{K3,E}(\PhiFA_3(D^b\Coh(K3\times E)))$}{Sp(K3,E) of PhiFA3}}
\label{spk3e:section}

Let
\[
 X=K3\times E,\qquad \pi_E:X\longrightarrow E,\qquad
 \cC_X=D^b\Coh(X),
\]
and let
\[
 \cF_X=\PhiFA_3(\cC_X)\in \Ethree\text{-}\mathrm{HolFA}(X)
\]
be the stage-one holomorphic factorisation algebra on the fixed framed
CY$_3$ branch. The pair \((K3,E)\) is a stage-two datum: the K3 factor
is the holomorphic surface integrated out and \(E\) is the reference
curve on which the residual chiral algebra lives.

\subsection{Relative Dolbeault Pushforward}
\label{spk3e:relative-dolbeault-pushforward}

\begin{definition}[K3-fibre chiral specialisation]
\label{spk3e:def-specialisation}
Assume that \(\cF_X\) is represented by a Costello-Gwilliam-Li
Dolbeault model of observables \((\Obs_X,Q+\dbar_X)\). For a Stein disk
\(U\subset E\), define
\[
 \SpCh_{K3,E}(\cF_X)(U)
 :=
 R\Gamma\!\left(K3\times U,\Obs_X\right)
 \underset{R\Gamma(K3,\Omega_{K3}^{0,\bullet})}{\otimes^{\mathbf L}}
 \CC_{\operatorname{tr}_{K3}},
\]
where \(\CC_{\operatorname{tr}_{K3}}\) is the one-dimensional module
defined by the Calabi-Yau trace
\[
 \operatorname{tr}_{K3}(\alpha)
 = \int_{K3}\alpha\wedge \Omega_{K3},
 \qquad
 \alpha\in \Omega^{0,2}(K3).
\]
The factorisation product is induced by disjoint disks in \(E\). Hence
\(\SpCh_{K3,E}(\cF_X)\) is an \(\Eone\)-chiral factorisation algebra on
\(E\).
\end{definition}

\begin{proof}
Harmonic projection on the compact K3 surface gives a quasi-isomorphism
\[
 R\Gamma(K3,\Omega_{K3}^{p,\bullet})
 \simeq
 \bigoplus_q H^q(K3,\Omega^p_{K3}).
\]
The trace above contracts the relative K3 Dolbeault direction and leaves
only the factorisation direction on \(U\subset E\). Dunn additivity for
the holomorphic splitting \(2+1=3\) identifies the residual
one-dimensional factorisation product with the \(\Eone\)-chiral product
on the curve. The braided \(\Etwo\)-structure at \(d\geq 3\) belongs to
\(\mathcal Z(\mathrm{Rep}^{\Eone}(-))\), equivalently to the derived
chiral centre, and is separate from the native algebra on \(E\).
\end{proof}

\subsection{Mukai-Heisenberg Zero Modes}
\label{spk3e:mukai-heisenberg-zero-modes}

\begin{proposition}[The abelian sector]
\label{spk3e:prop-abelian-sector}
Let \(\cF^{\mathrm{ab}}_X\subset \cF_X\) be the abelian holomorphic
BF sector with gauge Lie algebra \(\CC\). Then
\[
 \SpCh_{K3,E}\!\left(\cF^{\mathrm{ab}}_X\right)
 \simeq
 \Uchir(\widehat{\fh}_{\widetilde H(K3,\ZZ)})
\]
on \(E\). The lattice is the Mukai lattice
\[
 \widetilde H(K3,\ZZ)=H^0(K3,\ZZ)\oplus H^2(K3,\ZZ)\oplus H^4(K3,\ZZ)
 \cong \ZZ\oplus \mathrm{II}_{3,19}\oplus \ZZ
\]
of rank \(24\) and signature \((4,20)\), with pairing
\[
 \bigl((r,\xi,s),(r',\xi',s')\bigr)_{\Muk}
 = \xi\cdot \xi'-rs'-r's.
\]
Equivalently the generating currents satisfy
\[
 J_v(z_1)J_w(z_2)\sim
 \frac{(v,w)_{\Muk}}{(z_1-z_2)^2},
 \qquad v,w\in\widetilde H(K3,\ZZ).
\]
\end{proposition}

\begin{proof}
The abelian Dolbeault observables are a free commutative BV algebra on
the Dolbeault complex of \(X\). Relative Hodge projection gives
\[
 R\Gamma(K3\times U,\Omega_X^{\bullet,\bullet})
 \simeq
 H^\bullet(K3,\CC)\otimes R\Gamma(U,\Omega_E^{\bullet,\bullet}).
\]
The K3 Poincare pairing and the degree shift induced by the CY trace are
the Mukai pairing. Quantisation of a free chiral boson with this lattice
of zero modes is the chiral enveloping algebra of the affine Heisenberg
\(\widehat{\fh}_{\widetilde H(K3,\ZZ)}\). Restricting to \(H^2(K3,\ZZ)\)
recovers the K3 lattice \(\mathrm{II}_{3,19}\); the full Mukai extension
adds the two isotropic generators \(H^0\) and \(H^4\).
\end{proof}

\subsection{Finite Hall Positive Half}
\label{spk3e:finite-hall-positive-half}

\begin{definition}[Positive compact Hall windows]
\label{spk3e:def-hall-window}
Let \(\Gamma_+\) be the effective charge cone of reduced sheaf classes on
\(K3\times E\), and let \(\operatorname{ht}:\Gamma_+\to\ZZ_{\geq0}\) be
a height function. For \(h\geq0\), set
\[
 H^+_{\leq h}(X)
 :=
 \bigoplus_{\substack{\gamma\in\Gamma_+\\ \operatorname{ht}(\gamma)\leq h}}
 H^\bullet_{\BM,\mathrm{red}}\!\left(\cM_\gamma(X),\varphi_{\operatorname{tr}W}\right),
\]
where \(\cM_\gamma(X)\) is the derived moduli stack in class \(\gamma\)
and \(\varphi_{\operatorname{tr}W}\) is the critical CoHA perverse sheaf.
The Hall product is the correspondence product
\[
 m_{\gamma_1,\gamma_2}
 =
 (p_3)_!\,(p_1,p_2)^*
 :
 H^+_{\gamma_1}(X)\otimes H^+_{\gamma_2}(X)
 \longrightarrow H^+_{\gamma_1+\gamma_2}(X)
\]
through the stack of short exact sequences
\(\mathcal E_{\gamma_1,\gamma_2}(X)\).
\end{definition}

\begin{proposition}[What the BPS construction supplies]
\label{spk3e:prop-bps-positive}
The construction in Definition~\ref{spk3e:def-hall-window} supplies an
associative positive Hall half, filtered by height. Its primitive part
maps, after choosing the K3 reduced orientation and the Gritsenko-Nikulin
parity datum, to the positive nilpotent target
\[
 \operatorname{Prim}\,H^+_{\leq h}(X)
 \longrightarrow
 U\!\left(\fn_+(\fg_{\Delta_5})\right)_{\leq h}.
\]
This is a positive-half statement. It is distinct from a compact
Hall-Drinfeld double and from the Mukai Yangian branch.
\end{proposition}

\begin{proof}
Associativity is the standard Hall associativity of the two-step
extension stack. The target \(U(\fn_+(\fg_{\Delta_5}))\) is fixed by the
signed root-character data of the Gritsenko-Nikulin denominator:
for a positive root \(\alpha=(n,\ell,m)\),
\[
 \operatorname{sdim}\fg_{\Delta_5,\alpha}
 =
 f(nm,\ell),
\]
where \(f(n,\ell)\) is the Fourier coefficient of
\(\phi_{0,1}(\tau,z)\). A compact double requires a negative half, a
nondegenerate Hopf pairing, coproduct control, centre compatibility, and
completion. Those data are extra structure beyond the positive Hall
window.
\end{proof}

\subsection{Primitive Denominator and Scalar Square}
\label{spk3e:primitive-denominator}

\begin{theorem}[Primitive Borcherds denominator at \(N=1\)]
\label{spk3e:thm-delta5-denominator}
Let
\[
 \phi_{0,1}(\tau,z)=y^{-1}+10+y+O(q)
\]
be the K3 weak Jacobi form in Eichler-Zagier normalisation. The
Borcherds-Gritsenko-Nikulin lift gives the rank-three BKM Lie
superalgebra \(\fg_{\Delta_5}\) with denominator
\[
 D_5(2Z)
 =
 e^{2\pi i(\rho,Z)}
 \prod_{\alpha>0}
 \left(1-e^{2\pi i(\alpha,Z)}\right)^{
 \operatorname{sdim}\fg_{\Delta_5,\alpha}}
 =
 64^{-1}\Delta_5(2Z).
\]
Thus
\[
 D_5=64^{-1}\Delta_5,\qquad
 \Phi_{10}^{\mathrm{un}}=\Delta_5^2,\qquad
 \Phi_{10}^{\mathrm{OP}}=D_5^2=4096^{-1}\Delta_5^2.
\]
The BKM modular characteristic is
\[
 \kBKM(\Delta_5)
 =
 \operatorname{wt}(\Delta_5)
 =
 \frac{c_1(0)}{2}
 =
 \frac{10}{2}
 =
 5.
\]
\end{theorem}

\begin{proof}
Borcherds' singular-theta weight formula gives
\(\operatorname{wt}\BorLift(\phi)=c_\phi(0,0)/2\). For the K3 input
\(\phi_{0,1}\), the constant coefficient is \(c_1(0)=10\). The
Gritsenko-Nikulin product normalisation gives the monic product
\(D_5(2Z)=64^{-1}\Delta_5(2Z)\). Squaring removes the Maass sign character
and gives the scalar Igusa form; the primitive BKM weight remains attached
to \(\Delta_5\).
\end{proof}

\begin{corollary}[Primitive CHL constants]
\label{spk3e:cor-chl-constants}
For \(N\in\{1,2,3,4,6\}\), let \(F_N^{\mathrm{CHL}}\) be the primitive
CHL/BKM denominator attached to the \(g_N\)-twined K3 Jacobi input. Then
\[
 \bigl(c_N(0)\bigr)_{N=1,2,3,4,6}
 =
 (10,8,6,4,2),
 \qquad
 \bigl(\kBKM(F_N^{\mathrm{CHL}})\bigr)_{N=1,2,3,4,6}
 =
 (5,4,3,2,1).
\]
The scalar-square lift has weights \((10,8,6,4,2)\). At \(N=1\),
\[
 F_1^{\mathrm{CHL}}=\Delta_5,\qquad
 \Phi_{10}^{\mathrm{un}}=(F_1^{\mathrm{CHL}})^2.
\]
For \(N>1\), identification with a named Gritsenko-Clery catalogue row is
an additional row-map theorem.
\end{corollary}

\subsection{The Four \texorpdfstring{$\kappa_\bullet$}{subscripted invariant} Values}
\label{spk3e:four-invariant-values}

\begin{proposition}[Four constructions on \(K3\times E\)]
\label{spk3e:prop-four-invariants}
For \(Y=K3\times E\),
\[
\begin{array}{ccl}
 \kcat(Y) &=& \chi(\cO_Y)=\chi(\cO_{K3})\chi(\cO_E)=2\cdot0=0,\\[2mm]
 \kch(Y) &=& \displaystyle\sum_q(-1)^qh^{0,q}(Y)=1-1+1-1=0,\\[2mm]
 \kHeis(Y) &=& 3,\\[1mm]
 \kBKM(\Delta_5) &=& 5,\\[1mm]
 \kfib(Y) &=& \operatorname{rk}\widetilde H(K3,\ZZ)=24.
\end{array}
\]
The spectrum attached to the four named invariants is
\[
 \{\kcat,\kHeis,\kBKM,\kfib\}(Y)=\{0,3,5,24\}.
\]
The compact Hodge supertrace \(\kch(Y)=0\) is recorded separately from
the Heisenberg specialisation \(\kHeis(Y)=3\).
\end{proposition}

\begin{proof}
Kunneth gives \(\chi(\cO_{K3\times E})=2\cdot0=0\). The Hodge column of
\(K3\times E\) is \((h^{0,0},h^{0,1},h^{0,2},h^{0,3})=(1,1,1,1)\),
so the compact chiral Hodge supertrace is \(0\). The Heisenberg value
\(3\) is the stage-two Heisenberg-Mukai modular characteristic on the
curve \(E\), with \(2\) from K3 and \(1\) from the elliptic curve. The
BKM value \(5\) is Theorem~\ref{spk3e:thm-delta5-denominator}. The fibre
value \(24\) is the rank of the extended Mukai lattice.
\end{proof}

\subsection{Hall-Borcherds Recognition}
\label{spk3e:hall-borcherds-recognition}

\begin{conjecture}[Compact recognition theorem]
\label{spk3e:conj-recognition}
The compact specialisation satisfies
\[
 \SpCh_{K3,E}(\cF_X)
 \simeq
 \mathbf H_{\Delta_5}
\]
as an \(\Eone\)-chiral bialgebra on \(E\) after the following data have
been constructed:
\[
\begin{array}{ll}
 \mathrm{HB}_1 & \text{compact critical Hall source with K3 reduced orientation},\\
 \mathrm{HB}_2 & \text{negative Hall half and nondegenerate Hopf pairing},\\
 \mathrm{HB}_3 & \text{parity fixture matching }\operatorname{sdim}\fg_{\Delta_5,\alpha}=f(nm,\ell),\\
 \mathrm{HB}_4 & \text{coproduct, associator, and centre compatibility},\\
 \mathrm{HB}_5 & \text{completion exactness in charge height},\\
 \mathrm{HB}_6 & \text{denominator map } \operatorname{den}=D_5(2Z)=64^{-1}\Delta_5(2Z).
\end{array}
\]
\end{conjecture}

\begin{remark}[Two quantum-group branches]
\label{spk3e:two-branches}
The Mukai Yangian branch is the integrable quantisation of the rank-24
Mukai-Heisenberg current algebra. The BKM branch is the
Hall-Drinfeld/Borcherds branch controlled by
\(\fg_{\Delta_5}\) and its denominator \(D_5\). Standard Drinfeld
presentations apply to finite or Kac-Moody Cartan data. The algebra
\(\fg_{\Delta_5}\) has BKM imaginary simple roots, hence requires
Borcherds-Serre relations and signed root-character data. The comparison
target is therefore a completed
Hall-Drinfeld double of the compact Hall positive half, followed by the
Borcherds recognition map.
\end{remark}

\begin{remark}[Six constructions of \(G(K3\times E)\)]
\label{spk3e:six-constructions}
The six comparison routes to \(G(K3\times E)\) are six constructions:
Schiffmann-Vasserot surface CoHA, Maulik-Okounkov stable envelopes,
Gritsenko-Nikulin Borcherds denominator, Costello-Paquette holomorphic
Chern-Simons boundary, Bridgeland-Smith autoequivalence and line-operator
geometry, and Gopakumar-Vafa or reduced Donaldson-Thomas protected
indices. They are related by comparison maps to the same denominator
datum \(D_5\), with the compact recognition gates
\(\mathrm{HB}_1,\ldots,\mathrm{HB}_6\) supplying the missing common
source theorem.
\end{remark}
